
\section{同方向同频率简谐振动的合成}

\begin{theorem}
\begin{align*}
x &= x_1 + x_2 = A_1 \cos(\omega t + \phi_1) + A_2 \cos(\omega t + \phi_2)\\
&=A_1 \cos(\omega t)\cos\phi_1 - A_1 \sin(\omega t)\sin\phi_1+A_2 \cos(\omega t)\cos\phi_2 - A_2 \sin(\omega t)\sin\phi_2\\
&=\Big( A_1 \cos\phi_1 + A_2 \cos\phi_2 \Big) \cos(\omega t) - \Big( A_1 \sin\phi_1 + A_2 \sin\phi_2 \Big) \sin(\omega t)\\
&=\sqrt{A_1^2+ A_2^2 + 2A_1A_2(\cos\phi_1\cos\phi_2+\sin\phi_1\sin\phi_2)} \cos\left[ \omega t + \arctan\left( \frac{A_1 \sin\phi_1 + A_2 \sin\phi_2}{A_1 \cos\phi_1 + A_2 \cos\phi_2} \right) \right]\\
&= \sqrt{A_1^2 + A_2^2 + 2A_1A_2\cos(\phi_2 - \phi_1)} \cos(\omega t + \phi),~\tan\phi=\frac{A_1 \sin\phi_1 + A_2 \sin\phi_2}{A_1 \cos\phi_1 + A_2 \cos\phi_2}
\end{align*}
若两振动振幅相同，均为 \(A\)，且相位差为 \(\Delta\varphi\)，由合振幅公式与半角恒等式可得
 \[A_{\text{合}}^2=A^2+A^2+2A^2\cos\Delta\varphi =2A^2(1+\cos\Delta\varphi) =4A^2\cos^2\frac{\Delta\varphi}{2} \implies A_{\text{合}}=2A\cos\frac{\Delta\varphi}{2}.\]
所以同相时 \(\Delta\varphi=0\)，合振幅最大为 \(2A\)；反相时 \(\Delta\varphi=\pi\)，合振幅变成零。这正对应“完全相长”和“完全相消”。
\end{theorem}
本质上也是矢量的合成。比如一质点同时参与三个同方向、同频率简谐振动：
\[
x_1=A\cos\left(\omega t+\frac{\pi}{3}\right),\quad
x_2=A\cos\left(\omega t+\frac{5\pi}{3}\right),\quad
x_3=A\cos(\omega t+\pi).
\]
三者相隔 \(120^\circ\)，矢量相加后正好为 \(x=0\)。三个等长相量首尾相接，围成闭合三角形。
\begin{center}
\includegraphics[width=0.8\linewidth]{pdf-exercises/1.png}
\end{center}
